\documentclass[11pt]{article}
\usepackage[a4paper,margin=29mm]{geometry}
\usepackage[T1]{fontenc}
\usepackage{lmodern,amsmath,amssymb,amsthm,mathtools,microtype}
\usepackage[colorlinks=true,linkcolor=blue,citecolor=blue,urlcolor=blue]{hyperref}
\newtheorem{theorem}{Theorem}[section]
\newtheorem{proposition}[theorem]{Proposition}
\newtheorem{lemma}[theorem]{Lemma}
\newtheorem{remark}[theorem]{Remark}
\newcommand{\R}{\mathbb R}
\newcommand{\C}{\mathbb C}
\newcommand{\Tr}{\operatorname{Tr}}
\newcommand{\supp}{\operatorname{supp}}
\newcommand{\diver}{\operatorname{div}}
\newcommand{\Nc}{\mathcal N_c}
\newcommand{\N}{\mathcal N}
\newcommand{\HH}{\mathbb H}
\newcommand{\dd}{\,\mathrm d}
\newcommand{\norm}[1]{\lVert #1\rVert}
\newcommand{\abs}[1]{\lvert #1\rvert}
\title{Holomorphy of the Dirichlet-to-Neumann semigroup\newline
on continuous boundary functions for Lipschitz domains}
\author{Kaiwen Jin\thanks{Contact: \href{mailto:yui.ui.siki@gmail.com}{yui.ui.siki@gmail.com}.}}
\date{October 12, 2026}
\begin{document}
\maketitle
\begin{abstract}
Let $\Omega\subset\R^d$, $d\geq2$, be a bounded connected Lipschitz domain,
and let $C$ be a real symmetric uniformly elliptic matrix with Lipschitz
entries on $\overline\Omega$. We prove that the weak Dirichlet-to-Neumann
semigroup on $C(\partial\Omega)$ is holomorphic of angle $\pi/2$ and bounded
on every strictly smaller sector. This answers the continuous-function
question in the real symmetric case of Question 4 of ter Elst and Ouhabaz.
The proof obtains a supremum norm resolvent bound by boundary rescaling.
A complex Robin uniqueness lemma on a half-space excludes the limiting
solution; its proof uses vertical $L^2$ Dirichlet estimates and an iteration
of boundary mass bounds.
\end{abstract}

\section{Introduction}

The Dirichlet-to-Neumann operator describes diffusion on the boundary through
harmonic extension into the domain. Its realization on $L^2$ is nonnegative
and self-adjoint when the underlying elliptic coefficients are real and
symmetric. Holomorphy on continuous boundary functions requires an estimate
in the supremum norm, which does not follow from this Hilbert space fact.

Throughout, $\Omega\subset\R^d$, $d\geq2$, is a bounded connected Lipschitz
domain, $\Gamma=\partial\Omega$, and $\sigma$ is $(d-1)$-dimensional Hausdorff
measure on $\Gamma$. The matrix $C=(c_{kl})$ is real symmetric, its entries
are Lipschitz on $\overline\Omega$, and
\begin{equation}\label{eq:elliptic}
 \kappa\abs\xi^2\leq \xi\cdot C(x)\xi\leq K\abs\xi^2
 \qquad(x\in\overline\Omega,\ \xi\in\R^d)
\end{equation}
for some $0<\kappa\leq K<\infty$. All function spaces below are complex.
The weak Dirichlet-to-Neumann operator $\N$ on $L^2(\Gamma,\sigma)$ is defined
by $\varphi\in D(\N)$ and $\N\varphi=h$ if there is $u\in H^1(\Omega)$ with
$\Tr u=\varphi$ such that
\begin{equation}\label{eq:dtn}
 \int_\Omega C\nabla u\cdot\nabla\overline v\dd x
   =\int_\Gamma h\,\overline{\Tr v}\dd\sigma
 \qquad(v\in H^1(\Omega)).
\end{equation}
In particular, $u$ is harmonic for $-\diver(C\nabla)$, and is the unique
harmonic extension of $\varphi$. The sign in \eqref{eq:dtn} is that of the
outward conormal derivative. Write $\Nc$ for the part of $\N$ in
$C(\Gamma)$:
\[
 D(\Nc)=\{\varphi\in D(\N)\cap C(\Gamma):\N\varphi\in C(\Gamma)\}.
\]
We equip $C(\Gamma)$ with its usual supremum norm.

Arendt and ter Elst \cite{ArendtElst} proved that $-\Nc$ generates a positive
contractive $C_0$-semigroup $T$ and that $T(t)$ has a continuous kernel for
$t>0$. Their result identifies this semigroup with the restriction of the
$L^2$ semigroup. On $C^{1+\alpha}$ domains, ter Elst and Ouhabaz
\cite{ElstOuhabazAnalytic} established holomorphy of angle $\pi/2$ under
H\"older regularity of real symmetric coefficients. The final paragraph of
Question 4 in \cite{OpenProblems} asks about holomorphy on $C(\Gamma)$ when
the boundary is only Lipschitz and the real coefficients are Lipschitz.
We prove the following result for the symmetric class above.

For $0<\theta\leq\pi/2$, put
$\Sigma_\theta=\{z\in\C\setminus\{0\}:\abs{\arg z}<\theta\}$.

\begin{theorem}\label{thm:main}
Under \eqref{eq:elliptic} and the preceding hypotheses, $T$ extends to an
operator norm holomorphic semigroup on $\Sigma_{\pi/2}$. For every
$0<\theta<\pi/2$ there is a finite constant $M_\theta$, depending on
$\Omega,C,\theta$, such that
\begin{equation}\label{eq:sectorT}
 \sup_{z\in\Sigma_\theta}\norm{T(z)}_{C(\Gamma)\to C(\Gamma)}\leq M_\theta.
\end{equation}
For every $\varphi\in C(\Gamma)$, $T(z)\varphi\to\varphi$ as $z\to0$ in any
strictly smaller sector. No connectedness assumption on $\Gamma$ is needed.
\end{theorem}

For every $\varepsilon\in(0,\pi)$ there is a finite $M_\varepsilon$ such that
\begin{equation}\label{eq:mainres}
 \norm{\lambda(\lambda+\Nc)^{-1}}_{C(\Gamma)\to C(\Gamma)}
 \leq M_\varepsilon
 \qquad(\lambda\neq0,\ \abs{\arg\lambda}\leq\pi-\varepsilon).
\end{equation}
If this estimate fails at large $\abs\lambda$, rescaling at a point of
maximum boundary value gives a bounded solution of a complex Robin problem
on an unbounded Lipschitz graph. The original coefficient field becomes
constant under rescaling. Flattening the limiting graph then gives a
uniformly elliptic matrix independent of the normal variable. For that
matrix, the half-space $L^2$ Dirichlet theorem of
Auscher, Axelsson and McIntosh \cite{AAM} controls the mass of a bounded
harmonic function above a large boundary ball. Combining this estimate with
the Robin energy identity repeatedly improves the power of the boundary
mass until the limiting solution vanishes.

The graph area densities need not converge to the area density of the
limiting graph. We keep their weak-$*$ limit as a separate positive bounded
weight in the Robin condition. This is why the half-space uniqueness result
below includes such a weight.

\section{Half-space estimates and complex Robin uniqueness}

Let $n\geq1$, $\HH=\R^n\times(0,\infty)$, and write its coordinates as
$(x,t)$. In this section $A=A(x)$ is a bounded measurable real symmetric
$(n+1)\times(n+1)$ matrix, independent of $t$, satisfying uniform ellipticity
as in \eqref{eq:elliptic}. We use $B'(q,R)$ for a ball in $\R^n$ and
\[
 Q(q,R)=B'(q,R)\times(0,R).
\]
The phrase $H^1_{\mathrm{loc}}(\overline\HH)$ means $H^1$ on every bounded
half-space patch, including its flat boundary. Compactly supported tests in
this space are allowed to have a nonzero boundary trace.

\subsection{Dirichlet solutions of finite energy}

Let $\dot H^1_0(\HH)$ be the completion of $C_c^\infty(\HH)$ for the
gradient norm. Hardy's inequality
\begin{equation}\label{eq:Hardy}
 \norm{W/t}_{L^2(\HH)}\leq2\norm{\partial_t W}_{L^2(\HH)}
\end{equation}
identifies its elements as local Sobolev functions and fixes their additive
constant. A local Sobolev function with zero trace and finite gradient norm
belongs to this space. Indeed, vertical integration by parts on $(0,R)$,
with a nonpositive upper boundary term, proves \eqref{eq:Hardy}; radial
cutoffs and then cutoffs near $t=0$ have errors controlled by the tails of
$\abs{\nabla W}^2+\abs{W/t}^2$. Mollification of the remaining compact
interior functions gives the assertion.

For compactly supported $f\in H^{1/2}(\R^n)$, its usual Poisson extension
$P_t f$ has finite gradient norm. With the unitary Fourier transform,
\begin{equation}\label{eq:Poissonenergy}
 \int_\HH\abs{\nabla P_t f}^2\dd x\dd t
     =\int_{\R^n}\abs\xi\abs{\widehat f(\xi)}^2\dd\xi.
\end{equation}
This follows by integrating the multiplier $e^{-t\abs\xi}$ in $t$ and using
Plancherel \cite[Chapter 9]{RudinRCA}. The Poisson kernel is nonnegative,
has total mass one, and contracts $L^2$ and $L^\infty$.
Lax--Milgram on $\dot H^1_0(\HH)$ supplies a unique $A$-harmonic function
$H_Af$ such that $H_Af-P_tf\in\dot H^1_0(\HH)$. We use this notation only
for the variational solution just defined.

\begin{lemma}\label{lem:Dirichlet}
There is a finite constant $K_A$ such that, for every compactly supported
$f\in H^{1/2}(\R^n)$,
\begin{equation}\label{eq:vertical}
 \sup_{t>0}\norm{H_Af(\cdot,t)}_{L^2(\R^n)}\leq K_A\norm f_{L^2(\R^n)}.
\end{equation}
If also $f\in L^\infty$, then
\begin{equation}\label{eq:maxhalf}
 \norm{H_Af}_{L^\infty(\HH)}\leq\norm f_{L^\infty(\R^n)}.
\end{equation}
\end{lemma}

\begin{proof}
For smooth compactly supported $f$, the solution agrees with the $L^2$
Dirichlet solution of \cite{AAM}, by Theorem 5.3 there. Its
connected-component hypothesis holds because $(1-s)I+sA$, $0\leq s\leq1$,
stays in the real symmetric uniformly elliptic class; this class is
well posed by their Theorem 2.2(i). Their Corollary 2.4 gives
\eqref{eq:vertical}. The homogeneous zero-trace space in their Definition 5.1
agrees with the completion above by \eqref{eq:Hardy} and the cutoff argument.
Their ordering $(t,x)$ is changed to $(x,t)$ by an orthogonal permutation.

Approximate general $f$ by smooth compactly supported $f_k$ in $H^{1/2}$.
Formula \eqref{eq:Poissonenergy} and the variational equation give convergence
of the gradients. Hardy's inequality for the zero-trace differences then
gives local $H^1$ convergence up to $t=0$. Applied to differences, the
$L^2$ Dirichlet bound also gives convergence in $\sup_t L^2$ norm. These
limits agree locally, proving \eqref{eq:vertical} for the variational
solution.

For the maximum principle, first take real $f$ and set
$M=\norm f_\infty$. If $U=H_Af$, then
$w=(U-M)^+$ has finite gradient norm, zero trace, and
$w\leq\abs{U-P_tf}$, since $P_tf\leq M$. Thus $w/t\in L^2$, and the same
cutoff argument shows $w\in\dot H^1_0$. Testing the equation with $w$ gives
$\int A\nabla w\cdot\nabla w=0$, so $w=0$. The lower bound follows in the
same way. For complex data apply the upper bound to
$\operatorname{Re}(e^{-i\alpha}U)$ for a countable dense set of phases
$\alpha$. This yields \eqref{eq:maxhalf}.
\end{proof}

We will also need uniqueness without a global energy assumption.

\begin{lemma}\label{lem:zeroD}
A bounded $A$-harmonic function in $H^1_{\mathrm{loc}}(\overline\HH)$ with
zero trace is zero. The same assertion holds when the real uniformly
elliptic coefficients depend on both $x$ and $t$.
\end{lemma}

\begin{proof}
Reflect the function oddly across $t=0$ and the matrix by
$A_-(x,t)=JA(x,-t)J$, where $J=\operatorname{diag}(I_n,-1)$.
For a whole-space test $\eta$, the upper-half-space test
$\eta(x,t)-\eta(x,-t)$ has zero trace. It proves the weak equation for the
odd extension, whose matrix remains real uniformly elliptic.
The interior H\"older estimate for such equations
\cite[Corollary 5.5]{Trudinger} gives, for each fixed ball $B_r$ and arbitrarily
large $R$, an oscillation bound $C(r/R)^\alpha\norm U_\infty$.
Hence each real component of the extension is constant. Oddness makes the
constant zero.
\end{proof}

\subsection{From boundary mass to bulk mass}

\begin{lemma}\label{lem:mass}
Let $U\in H^1_{\mathrm{loc}}(\overline\HH)$ be bounded and $A$-harmonic, and
write $\varphi=\Tr U$. Suppose that, for some $\beta<n$ and finite $M_\beta$,
\begin{equation}\label{eq:Fbeta}
 F(R):=\sup_{q\in\R^n}\int_{B'(q,R)}\abs\varphi^2\dd x
       \leq M_\beta R^\beta\qquad(R\geq1).
\end{equation}
Then there is a finite constant $C_\beta$ such that
\begin{equation}\label{eq:bulk}
 \sup_{q\in\R^n}\int_{Q(q,R)}\abs U^2\dd x\dd t
       \leq C_\beta R^{\beta+1}\qquad(R\geq1).
\end{equation}
\end{lemma}

\begin{proof}
Fix $q,R$ and a nonincreasing radial Lipschitz cutoff $\chi$, equal to one
on $B'_1$ and zero outside $B'_2$. Set
\[
 \chi_0(x)=\chi\left(\frac{x-q}{4R}\right),\qquad
 \chi_j(x)=\chi\left(\frac{x-q}{2^j4R}\right)
             -\chi\left(\frac{x-q}{2^{j-1}4R}\right)\quad(j\geq1).
\]
The $\chi_j$ are nonnegative and sum to one. The far piece $\chi_j$ vanishes
on $B'(q,S_j)$, where $S_j=2^{j-1}4R$. Put
$f_j=\chi_j\varphi$ and $U_j=H_Af_j$.
Every $f_j$ is compactly supported $H^{1/2}$ data: apply the local $H^1$
trace theorem to a tangential and vertical cutoff of $U$.

For the near piece, Lemma~\ref{lem:Dirichlet} gives
\begin{equation}\label{eq:near}
 \int_{Q(q,R)}\abs{U_0}^2
       \leq R K_A^2\norm{f_0}_2^2\leq C R^{\beta+1}.
\end{equation}
For $j\geq1$, odd reflection of $U_j$ gives a homogeneous elliptic equation
in the ball of radius $S_j$ centered at $(q,0)$. The interior
$L^2$-to-supremum estimate \cite[Theorem 5.1]{Trudinger}, together with
\eqref{eq:vertical}, implies
\begin{align}\label{eq:far}
 \sup_{Q(q,R)}\abs{U_j}
 &\leq C S_j^{-(n+1)/2}
      \norm{U_j}_{L^2(B'(q,S_j)\times(0,S_j))}\notag\\
 &\leq C S_j^{-n/2}\norm{f_j}_2
 \leq C_\beta S_j^{(\beta-n)/2}.
\end{align}
Here $Q(q,R)$ lies inside the half-radius ball because $S_j\geq4R$.
Since $\beta<n$, the geometric series in \eqref{eq:far} sums to at most
$C_\beta R^{(\beta-n)/2}$.

The sum $V=\sum_{j\geq0}U_j$ converges uniformly on compact patches for its
far tail. It also converges in local $H^1$ up to the boundary. To see the
latter assertion, fix a larger patch; sufficiently far finite tails have
zero trace there. Testing their equation with the tail times a squared
cutoff gives the boundary Caccioppoli estimate
\[
 \norm{\nabla(\text{tail})}_{L^2(\text{smaller patch})}^2
     \leq C\norm{\text{tail}}_{L^2(\text{larger patch})}^2.
\]
The uniform tail bound makes the right-hand side tend to zero. Thus $V$ is
harmonic with trace $\varphi$. By linearity its partial sums equal
$H_A((\sum_{j\leq k}\chi_j)\varphi)$, and \eqref{eq:maxhalf} bounds every
partial sum by $\norm\varphi_\infty$. Therefore $V$ is bounded and
Lemma~\ref{lem:zeroD} identifies it with $U$.
Combining \eqref{eq:near} with the square of the summed far bound and the
volume $O(R^{n+1})$ proves \eqref{eq:bulk}.
\end{proof}

\subsection{A Robin Liouville lemma}

\begin{proposition}\label{prop:Robin}
Let $b\in L^\infty(\R^n)$ satisfy $0<b_*\leq b\leq b^*<\infty$.
If $\lambda\in\C\setminus(-\infty,0]$ and a bounded
$U\in H^1_{\mathrm{loc}}(\overline\HH)$ satisfies
\begin{equation}\label{eq:Robin}
 \int_\HH A\nabla U\cdot\nabla\overline V\dd x\dd t
   +\lambda\int_{\R^n}b\,\varphi\,\overline{\Tr V}\dd x=0,
 \qquad \varphi=\Tr U,
\end{equation}
for all compactly supported local $H^1$ tests, then $U=0$.
\end{proposition}

\begin{proof}
Write $\alpha=\arg\lambda\in(-\pi,\pi)$. Choose a cutoff $\eta_R$ equal to
one on $B'(q,R)\times[0,R]$, supported in
$B'(q,2R)\times[0,2R]$, with $\abs{\nabla\eta_R}\leq C/R$.
Use $V=\eta_R^2 U$ in \eqref{eq:Robin}, multiply by $e^{-i\alpha/2}$,
and take real parts. Both the gradient quadratic form and the boundary
term are coercive, with the common positive factor $\cos(\alpha/2)$.
Young's inequality absorbs the cross term, giving
\begin{equation}\label{eq:Robinmass}
 \int_{B'(q,R)}\abs\varphi^2\dd x
    \leq C_\lambda R^{-2}\int_{Q(q,2R)}\abs U^2\dd x\dd t.
\end{equation}
The constants are finite for the fixed $A,b,\lambda$.

Boundedness of $U$ first gives $F(R)\leq C R^{n-1}$ for $R\geq1$.
This exponent is strictly below $n$. Whenever
$F(R)\leq C_\beta R^\beta$ with $\beta<n$, Lemma~\ref{lem:mass} and
\eqref{eq:Robinmass} improve it to $F(R)\leq C_{\beta-1}R^{\beta-1}$.
After $n$ improvements we obtain $F(R)\leq C/R$.
For any fixed ball, its trace mass is bounded by $F(R)$ for all sufficiently
large $R$, so it is zero. Hence $\varphi=0$, and Lemma~\ref{lem:zeroD}
gives $U=0$.
\end{proof}

\begin{remark}
The parameter exclusions in Proposition~\ref{prop:Robin} are necessary.
At $\lambda=0$ constants are allowed. If $A=I$, $b=1$, and
$\lambda=-\abs\xi<0$, then
$U(x,t)=e^{i\xi\cdot x}e^{-\abs\xi t}$ is a nonzero bounded solution.
\end{remark}

\section{Boundary compactness and the resolvent estimate}

\subsection{Local Neumann estimates on Lipschitz graphs}

The interior estimates used above are the uniformly elliptic specializations
of \cite[Theorem 5.1 and Corollaries 5.4--5.5]{Trudinger}.
For the inhomogeneous equation $\diver(B\nabla w)=\diver F$, bounded $F$
gives local H\"older continuity. After rescaling to fixed balls, the
constants depend only on dimension, ellipticity and the forcing bound.
The results apply separately to the real and imaginary parts.

\begin{lemma}\label{lem:Neumann}
Consider a Lipschitz graph $D=\{(x,s):s>\psi(x)\}$ with
$\operatorname{Lip}\psi\leq L$, a real uniformly elliptic matrix field $B_0$
with fixed bounds, and a harmonic $v\in H^1_{\mathrm{loc}}(\overline D)$
satisfying $\norm v_\infty\leq1$. Suppose that its weak outward conormal
derivative $h$ is bounded by $H$ on the graph patch in question.
On each smaller patch, $v$ has local $H^1$ and H\"older bounds depending only
on $d,L$, the ellipticity bounds, $H$ and the patch radii. It has ambient
extensions with the same local bounds across the graph.
\end{lemma}

\begin{proof}
Test the weak equation with $\eta^2v$. The boundary term is bounded by
$H$ times the surface area of the patch, and Young's inequality absorbs the
gradient cross term. This proves the local energy bound.

Flatten by $\Phi(x,t)=(x,t+\psi(x))$. The Sobolev chain rule and graph
reflection are valid for Lipschitz $\psi$;
see \cite[arXiv version, Section 3.1 and Lemma 3.3]{Nittka}. The determinant of $D\Phi$ is one,
and the transformed matrix is
\[
 B=(D\Phi)^{-1}B_0(\Phi)(D\Phi)^{-T}.
\]
It is real uniformly elliptic with bounds controlled by $L$.
The graph area formula \cite[Section 3.3.4]{EvansGariepy} gives
$a(x)=\sqrt{1+\abs{\nabla\psi(x)}^2}$.
Writing $\widetilde h(x)=h(x,\psi(x))$, the bounded field
$F=(0,\ldots,0,-\widetilde h a)$, independent of $t$, satisfies
\[
 \int_\HH F\cdot\nabla\overline V
     =\int_{\R^{d-1}}\widetilde h a\,\overline{\Tr V}.
\]
Thus the flattened weak equation is
$\int B\nabla(v\circ\Phi)\cdot\nabla\overline V=\int F\cdot\nabla\overline V$.
Extend $v\circ\Phi$ evenly across $t=0$, the matrix by $JBJ$, and the field
by $JF$, where $J=\operatorname{diag}(I_{d-1},-1)$. Testing with the
sum of the reflected tests proves the
whole-space inhomogeneous equation. The interior H\"older estimate now
applies. Transforming back and using the same reflection gives the claimed
ambient $H^1$ and H\"older extensions.
\end{proof}

\subsection{Compact resolvents and small parameters}

We first recall why all the resolvents used in \eqref{eq:mainres} exist and
why only large parameters can violate that estimate. The continuous kernel
theorem and the $C_0$ theorem in \cite[Corollary 3.5 and Theorem 5.3(a)]{ArendtElst}
apply with potential $V=0$: the coefficients are real symmetric and
Lipschitz, and the Dirichlet realization has no zero eigenvalue by
Poincar\'e's inequality. Positivity and contractivity follow from their
Proposition 5.1(b).

\begin{lemma}\label{lem:compactres}
For every $\lambda\in\C\setminus(-\infty,0]$, $\lambda+\Nc$ is invertible
on $C(\Gamma)$. Its resolvent is compact. For each $\varepsilon>0$,
$\lambda(\lambda+\Nc)^{-1}$ is uniformly bounded when
$0<\abs\lambda\leq R$ and $\abs{\arg\lambda}\leq\pi-\varepsilon$, for any
fixed finite $R$.
\end{lemma}

\begin{proof}
For $a>0$, contractivity and the semigroup identity give
\begin{equation}\label{eq:Laplace}
 R_a:=(a+\Nc)^{-1}=\int_0^\infty e^{-at}T(t)\dd t.
\end{equation}
For example, the difference quotient of the right-hand side $Q_af$ is
\[
 \frac{T(h)Q_af-Q_af}{h}
 =\frac{e^{ah}-1}{h}Q_af-\frac{e^{ah}}h\int_0^h e^{-as}T(s)f\dd s
 \longrightarrow aQ_af-f.
\]
It gives $Q_af\in D(\Nc)$ and $(a+\Nc)Q_af=f$. Injectivity follows from
the energy identity \eqref{eq:dtn}.

Every $T(t)$, $t>0$, is compact by its continuous kernel and Ascoli.
Strong continuity is uniform on the compact image of $T(t/2)$, so $T$ is
norm continuous for positive times. Truncating \eqref{eq:Laplace} to
$[\delta,M]$ and using norm Riemann sums proves that $R_1$ is compact.
For $\lambda\notin(-\infty,0]$, $I+(\lambda-1)R_1$ is injective: a kernel
vector belongs to $D(\Nc)$ and satisfies $(\lambda+\Nc)\varphi=0$, whereas
\begin{equation}\label{eq:energy}
 \lambda\norm\varphi_2^2+\int_\Omega C\nabla u\cdot\nabla\overline u=0
\end{equation}
forces $\varphi=0$ at such a parameter.
The compact Fredholm alternative therefore gives
\begin{equation}\label{eq:resformula}
 (\lambda+\Nc)^{-1}=R_1[I+(\lambda-1)R_1]^{-1}.
\end{equation}
One elementary form of that alternative is as follows. An injective $I-K$
with compact $K$ is bounded below, by a compactness argument on a sequence
of unit vectors. Its iterated ranges are closed. If its range were proper,
these ranges would strictly decrease; unit vectors separated from the next
range would make their images under $K$ uniformly separated, contradicting
compactness.

Let $P\varphi=\sigma(\Gamma)^{-1}\int_\Gamma\varphi\dd\sigma$, regarded as a
constant function. Connectedness of $\Omega$ and \eqref{eq:energy} give
$\ker\Nc=\C1$. The self-adjoint $L^2$ realization fixes constants and
preserves the boundary integral; hence $R_1$ commutes with $P$.
On $X_0=\ker P$, the compact Fredholm alternative makes $I-R_1$ invertible.
Since $\Nc R_1=I-R_1$, the restriction of $\Nc$ to $X_0$ has bounded inverse
$R_1(I-R_1)^{-1}$. The Neumann series at zero therefore bounds
$\lambda(\lambda+\Nc)^{-1}$ on $X_0$ for small $\abs\lambda$; on constants
it equals $P$. On the remaining compact parameter set, the resolvent
identity and \eqref{eq:resformula} give norm continuity and a finite bound.
\end{proof}

\subsection{The large-parameter estimate}

\begin{proposition}\label{prop:resolvent}
For every $\varepsilon\in(0,\pi)$ there is a finite $M_\varepsilon$ such that
\[
 \norm{\lambda(\lambda+\Nc)^{-1}}_{C(\Gamma)\to C(\Gamma)}
    \leq M_\varepsilon
 \qquad(\lambda\ne0,\ \abs{\arg\lambda}\leq\pi-\varepsilon),
\]
where $\Nc$ is the part of the weak Dirichlet-to-Neumann operator in $C(\Gamma)$.
\end{proposition}

\begin{proof}
By Lemma~\ref{lem:compactres}, a failure would give parameters
$\abs{\lambda_j}\to\infty$ with $\abs{\arg\lambda_j}\leq\pi-\varepsilon$,
data $f_j\in C(\Gamma)$, and
$\varphi_j=(\lambda_j+\Nc)^{-1}f_j$ such that
\begin{equation}\label{eq:fail}
 \frac{\norm{f_j}_\infty}{\abs{\lambda_j}\norm{\varphi_j}_\infty}\longrightarrow0.
\end{equation}
Choose $q_j\in\Gamma$ attaining the supremum and set
$\tau_j=\overline{\varphi_j(q_j)}/\abs{\varphi_j(q_j)}$. Let $u_j$ be the
harmonic extension of $\tau_j\varphi_j/\norm{\varphi_j}_\infty$, so that
$u_j(q_j)=1$. The weak maximum principle,
applied to real phase components, gives $\norm{u_j}_\infty\leq1$.
Its proof on this bounded domain is the usual test with the positive part
above the maximum boundary value, which has zero trace.

Pass to one Lipschitz chart containing the $q_j$ with a fixed margin, and
rotate coordinates once so that the domain is above its graph. Set
\[
 r_j=\abs{\lambda_j}^{-1},\qquad
 v_j(y)=u_j(q_j+r_jy),\qquad z_j=r_j\lambda_j.
\]
On every fixed bounded set, for large $j$ the scaled domain is
$D_j=\{y_d>\psi_j(y')\}$, where $\psi_j(0)=0$ and the Lipschitz constants
are bounded by a common $L$. A diagonal subsequence converges locally
uniformly to a global Lipschitz function $\psi$.
The coefficients $C_j(y)=C(q_j+r_jy)$, in the fixed rotated coordinates,
converge uniformly locally on $D_j$ to a constant positive matrix $C_*$.
Indeed, $q_j$ has a convergent subsequence and the Lipschitz modulus of $C$
at scale $r_j$ tends to zero.

The scaled weak equation is
\begin{equation}\label{eq:scaled}
 \int_{D_j} C_j\nabla v_j\cdot\nabla\overline\eta
   +z_j\int_{\partial D_j}v_j\overline\eta\dd\sigma_j
     =\int_{\partial D_j}g_j\overline\eta\dd\sigma_j,
\end{equation}
for compactly supported tests in the chart, where
\[
 g_j(y)=\frac{r_j\tau_j f_j(q_j+r_jy)}{\norm{\varphi_j}_\infty}
 \quad (y\in\partial D_j).
\]
Thus $\norm{g_j}_\infty\to0$ by \eqref{eq:fail}. The factor $r_j$ in the
boundary terms comes from dividing the weak equation by $r_j^{d-2}$,
while surface measure scales by $r_j^{d-1}$.
After a subsequence $z_j\to z_*$, where $\abs{z_*}=1$ and
$\abs{\arg z_*}\leq\pi-\varepsilon$.
The conormal derivatives $g_j-z_jv_j$ are uniformly bounded.
Lemma~\ref{lem:Neumann} gives reflected ambient extensions with uniform
local $H^1$ and H\"older bounds. Weak $H^1$ compactness and Ascoli on
successive compact sets yield a bounded limit $v$ on
$D=\{y_d>\psi(y')\}$, locally $H^1$ up to its boundary, with $v(0)=1$.

In graph coordinates set
$a_j(x)=\sqrt{1+\abs{\nabla\psi_j(x)}^2}$.
A further diagonal weak-$*$ subsequence gives
\begin{equation}\label{eq:density}
 a_j\mathrel{\overset{*}{\rightharpoonup}} b\quad\text{locally in }L^\infty,
 \qquad 1\leq b\leq\sqrt{1+L^2}.
\end{equation}
Uniform convergence of the extensions and of $\psi_j$ gives uniform
convergence of $v_j(x,\psi_j(x))$ on compact sets. These values can therefore
multiply the weak-$*$ convergence in \eqref{eq:density}.
In the bulk, $1_{D_j}\to1_D$ almost everywhere and strongly in every finite
local $L^p$ space. Together with coefficient convergence and weak convergence
of the extension gradients, this passes the bulk integral in
\eqref{eq:scaled}. The result is
\begin{equation}\label{eq:limit}
 \int_D C_*\nabla v\cdot\nabla\overline\eta
   +z_*\int_{\R^{d-1}}b(x)v(x,\psi(x))
                            \overline{\eta(x,\psi(x))}\dd x=0.
\end{equation}
Initially this holds for compact smooth ambient tests. They are dense in
the local graph $H^1$ test space: reflect across the graph, multiply by an
ambient cutoff, mollify the extension, and restrict back to $D$.
Trace continuity passes the boundary term as well.

Flatten the limiting graph by $\Phi(x,t)=(x,t+\psi(x))$. Its determinant is
one, so the matrix in the half-space equation is
\[
 A(x)=(D\Phi)^{-1}C_*(D\Phi)^{-T}.
\]
It is real symmetric uniformly elliptic and independent of $t$.
Equation \eqref{eq:limit} becomes precisely \eqref{eq:Robin} with
$U=v\circ\Phi$, $\lambda=z_*$ and the weight $b$ in \eqref{eq:density}.
Proposition~\ref{prop:Robin} gives $U=0$, contradicting $v(0)=1$.
This proves the large-parameter estimate and hence \eqref{eq:mainres}.
\end{proof}

\section{Holomorphic extension}

We finish with the contour argument, including the behavior at the vertex
of the sector. Put $X=C(\Gamma)$. With $\mu=-\lambda$,
Proposition~\ref{prop:resolvent} says that, for each $\omega>0$,
\begin{equation}\label{eq:sectorres}
 R(\mu)=(\mu-\Nc)^{-1},\qquad
 \norm{R(\mu)}_{X\to X}\leq\frac{M_\omega}{\abs\mu}
 \quad(\mu\neq0,\ \abs{\arg\mu}\geq\omega).
\end{equation}

\begin{proof}[Proof of Theorem~\ref{thm:main}]
Fix $0<\omega<\pi/2$ and $r>0$. Let $G_{\omega,r}$ run inward from infinity
on the ray $\arg\mu=\omega$, follow the circle of radius $r$
counterclockwise from angle $\omega$ to $2\pi-\omega$ through the negative
half-plane, and then run outward on the ray $\arg\mu=-\omega$.
For $\abs{\arg z}<\pi/2-\omega$, define
\begin{equation}\label{eq:contour}
 E(z)=\frac1{2\pi i}\int_{G_{\omega,r}}e^{-z\mu}R(\mu)\dd\mu.
\end{equation}
The integral converges in operator norm. Cauchy's theorem in the annular
region outside the positive sector makes it independent of $r$.
Banach-valued integration and Cauchy's theorem can be applied in norm here;
see \cite[Theorems 3.27, 3.29 and 3.31]{RudinFA}.

Fix $\theta<\pi/2-\omega$ and estimate with $r=1/\abs z$.
On either ray $\operatorname{Re}(z\mu)\geq c_\theta\abs z\abs\mu$.
Each ray contribution is at most
$M_\omega\int_1^\infty e^{-c_\theta s}\dd s/s$.
On the circle, $\abs{e^{-z\mu}}\leq e$ and its length is
$(2\pi-2\omega)r$. Thus $E$ is bounded for all $z\in\Sigma_\theta$.
Keeping $r$ fixed locally in $z$, differentiation under the integral proves
operator norm holomorphy.

For real $a>r$, absolute convergence permits integrating
$e^{-at}E(t)$ in $t$ first. We obtain
\[
 \int_0^\infty e^{-at}E(t)\dd t
   =\frac1{2\pi i}\int_{G_{\omega,r}}(a+\mu)^{-1}R(\mu)\dd\mu
   =(a+\Nc)^{-1}.
\]
For the last equality close the contour on the left. Its orientation is
clockwise, and the pole at $\mu=-a$ has residue
$R(-a)=-(a+\Nc)^{-1}$; the outer arc tends to zero by
\eqref{eq:sectorres}. This agrees with the Laplace transform of $T(t)$ in
\eqref{eq:Laplace}, so uniqueness gives $E(t)=T(t)$ for $t>0$.
For completeness, scalar uniqueness for bounded continuous functions follows
by the substitution $s=e^{-t}$: vanishing Laplace transforms give all
sufficiently high integer moments on $(0,1)$. Polynomial density in
$C[0,1]$ then gives zero against continuous tests and hence zero everywhere
by continuity on $(0,1)$. Bounded linear functionals separate the
Banach-valued functions.

The identity theorem extends the semigroup law from positive real parameters
to $E(z+w)=E(z)E(w)$ on each such sector, first in one variable and then in
the other. The sector is convex and closed under addition. Constructions
with different $\omega$ agree on overlaps, and their union is
$\Sigma_{\pi/2}$. Denote the resulting extension by $T(z)$.

Finally, for $\varphi\in D(\Nc)$ the sector bound makes
$\int_0^z T(w)\Nc\varphi\dd w$ a holomorphic primitive from the vertex.
One may integrate first from an interior point and use Cauchy's theorem;
small connecting arcs at zero have norm $O(\text{radius})$, so the limit is
path independent. On positive real parameters the semigroup derivative
gives this primitive as $\varphi-T(z)\varphi$. The identity theorem gives
the same equality throughout the sector, and hence
\[
 \norm{T(z)\varphi-\varphi}_\infty
    \leq M_\theta\abs z\norm{\Nc\varphi}_\infty
       \qquad(z\in\Sigma_\theta).
\]
Density of $D(\Nc)$ in $X$, together with \eqref{eq:sectorT}, proves strong
continuity at zero for every $\varphi\in X$.
\end{proof}

\paragraph{AI assistance.}
OpenAI Codex assisted with proof exploration, checking source statements,
and preparing the manuscript. The proof has undergone an AI natural-language
review; it has not been formally verified or independently certified by a
human referee.

\begin{thebibliography}{10}
\bibitem{ArendtElst}
W. Arendt and A. F. M. ter Elst,
\emph{The Dirichlet-to-Neumann operator on $C(\partial\Omega)$},
Ann. Sc. Norm. Super. Pisa Cl. Sci. (5) \textbf{20} (2020), 1169--1196.
\href{https://arxiv.org/abs/1707.05556}{arXiv:1707.05556}.

\bibitem{AAM}
P. Auscher, A. Axelsson and A. McIntosh,
\emph{Solvability of elliptic systems with square integrable boundary data},
Ark. Mat. \textbf{48} (2010), 253--287.
\href{https://doi.org/10.1007/s11512-009-0108-2}{doi:10.1007/s11512-009-0108-2}.

\bibitem{ElstOuhabazAnalytic}
A. F. M. ter Elst and E. M. Ouhabaz,
\emph{Analyticity of the Dirichlet-to-Neumann semigroup on continuous functions},
J. Evol. Equ. \textbf{19} (2019), 21--31.
\href{https://arxiv.org/abs/1707.07718}{arXiv:1707.07718}.

\bibitem{OpenProblems}
A. F. M. ter Elst and E. M. Ouhabaz,
\emph{Five open problems on the Dirichlet-to-Neumann operator with variable
coefficients}, Integral Equations Operator Theory \textbf{97} (2025), article 14.
\href{https://doi.org/10.1007/s00020-025-02796-9}{doi:10.1007/s00020-025-02796-9}.

\bibitem{EvansGariepy}
L. C. Evans and R. F. Gariepy,
\emph{Measure Theory and Fine Properties of Functions}, revised edition,
CRC Press, 2015.

\bibitem{Nittka}
R. Nittka,
\emph{Regularity of solutions of linear second order elliptic and parabolic
boundary value problems on Lipschitz domains},
J. Differential Equations \textbf{251} (2011), 860--880.
\href{https://doi.org/10.1016/j.jde.2011.05.019}{doi:10.1016/j.jde.2011.05.019}.
Preprint: \href{https://arxiv.org/abs/0906.5285v1}{arXiv:0906.5285v1}.

\bibitem{RudinRCA}
W. Rudin, \emph{Real and Complex Analysis}, third edition,
McGraw--Hill, 1987.

\bibitem{RudinFA}
W. Rudin, \emph{Functional Analysis}, second edition,
McGraw--Hill, 1991.

\bibitem{Trudinger}
N. S. Trudinger,
\emph{Linear elliptic operators with measurable coefficients},
Ann. Scuola Norm. Sup. Pisa Cl. Sci. (3) \textbf{27} (1973), 265--308.
\href{https://www.numdam.org/item/ASNSP_1973_3_27_2_265_0/}{Numdam}.
\end{thebibliography}
\end{document}
